\chapter{O que falta}
\label{cap:falta}

\section{Perguntas em aberto}

\campo{Alocação do corte entre usinas} A seção~\ref{sec:fig10} identificou duas
covariáveis que sobrevivem à correção de multiplicidade: distância de rede à
demanda e capacidade da usina. Sob os critérios da seção~\ref{sec:fig12} é a
segunda que resiste, e é a que não tem explicação testada. A variância explicada
é pequena, e a maior parte da dispersão da taxa de corte permanece sem
explicação.

\campo{Causa do corte no Nordeste} Nem piso térmico local, nem escoamento. Resta
ausência de demanda no destino, obtida por eliminação e não testada de forma
positiva.

\campo{Inflexibilidade declarada contra mínimo técnico real} É a grandeza que
uma auditoria exigiria, e não está publicada.

\section{Trabalho seguinte, em ordem de retorno esperado}

\begin{enumerate}[leftmargin=1.2em,itemsep=0.35em]
\item \textbf{Reconstruir o corte anterior a abril de 2024.} A maquinaria da
seção~\ref{sec:fig6} já existe. Um modelo treinado nos intervalos sem restrição,
aplicado de forma retroativa, viabiliza o teste temporal da
seção~\ref{sec:fig4}. Há verdade de campo para validar, nos dois anos em que as
duas séries coexistem.

\item \textbf{Medir o resíduo não explicado.} No nível de intervalo, a diferença
de desempenho entre um modelo só com observáveis sistêmicos e outro com
identidade de usina quantifica o que os capítulos de rede deixaram sem
explicação. A série eólica acrescenta 1,4 GB e três anos de painel. A primeira
metade já está publicada: Gorka e Roald \cite{gorka2024} classificam a
ocorrência de corte no CAISO a partir do estado do sistema, que é o modelo sem
identidade. O que não se encontrou publicado é o passo seguinte, de recalcular o
desempenho \emph{dentro} de cada intervalo, onde a informação de quando o corte
ocorre é removida por construção e resta a de qual usina o recebe.

\item \textbf{Substituir NASA POWER por NSRDB.} A célula passa de cerca de 50 km
para 2 km, e o passo de diário para 10 minutos, sem custo. Atende três das
objeções já escritas.

\item \textbf{Delinear o arranjo por segmentação.} Substitui os círculos de
400 m da seção~\ref{sec:fig7}, aumenta a precisão da datação e mede a razão
entre terreno e módulo que o capítulo~\ref{cap:aplicacao} assume.
\end{enumerate}

\section{Uma regra de método}

Qualquer saída de modelo que entre em inferência a jusante enviesa o ponto e o
erro padrão \cite{proctor2023}. A predição entra como distribuição, e não como
valor pontual. A regra vale para a irradiância do capítulo~\ref{cap:calibracao}
e para o grafo do capítulo~\ref{cap:rede}, e não foi aplicada em nenhum dos
dois.

\campo{O cálculo está disponível} Angelopoulos et al. \cite{angelopoulos2023}
dão o procedimento: intervalo de confiança válido para média, quantil e
coeficiente de regressão, sem hipótese sobre o algoritmo que produz a predição,
desde que exista um subconjunto com o valor observado. Kluger et al.
\cite{kluger2025} estendem o procedimento ao caso em que a predição entra como
covariável imputada, e motivam com predição de indicador ambiental por imagem de
satélite alimentando regressão a jusante, que é a forma exata dos capítulos de
calibração e de rede.

A condição é ter verdade de campo num subconjunto. Para a irradiância ela
existe: a série medida no plano do arranjo cobre a mesma janela. Para o grafo
não há equivalente publicado, e ali a correção exige antes definir contra o quê
a predição seria comparada.

\begin{objecao}
Aplicá-la alargaria os intervalos de todos os resultados de rede. Nenhum deles
foi medido sob essa correção, e não se pode afirmar que o resultado que passa
nos três critérios sobreviveria a ela.
\end{objecao}
